%% kapa: exam=1 round=28 year=2017 period=2 form=A subject=accounting subject_ko=회계학 seq=12 number=52
%% source: Q-Net 감정평가사 기출문제 (cst00309) artlSeq=5199603 file=extracted/1008/5199603/2017년 제28회 감정평가사 1차시험 문제지/2교시 A형.pdf
%% conversion: kordoc 4.24.51 -> tools/gen1.py (auto); review=proofread-sonnet; flags=none
\begin{question}{52}
(주)대한은 20x1년 1월 1일에 건물을 ￦20,000,000에 취득하여 사용하고 있다 (내용연수 5년, 잔존가치 ￦0, 정액법 상각). (주)대한은 20x2년 말에 재평가모형을 최초 적용하였으며, 장부금액과 감가상각누계액을 비례하여 수정하는 방법으로 회계처리하고 있다. 20x2년 말 건물의 재평가 전 장부금액은 ￦12,000,000이고, 공정가치는 ￦18,000,000이다. 다음 중 옳지 않은 것은? (단, 재평가잉여금은 이익잉여금으로 대체하지 않는다.)
\begin{choices}
  \item (주)대한은 20x2년에 감가상각비로 ￦4,000,000을 보고한다.
  \item 재평가 이후 (주)대한의 20x2년 말 재무상태표에 표시되는 건물의 장부금액은 ￦18,000,000이다.
  \item (주)대한은 20x2년 말에 재무상태표에 ￦6,000,000을 기타포괄손익누계액으로 인식한다.
  \item 자산재평가로 인한 장부금액 수정을 기존의 감가상각누계액을 제거하는 방법을 사용하면 건물의 장부금액 ￦18,000,000이 보고된다.
  \item (주)대한이 20x1년에도 이 건물에 대한 재평가를 실시하여 재평가손실 ￦1,000,000을 인식하였다면, 20x2년에는 ￦5,000,000을 당기이익으로 인식한다.
\end{choices}
\end{question}
\answer{5}
